\section{Accessing nested fields}\label{sec:08-rings:06-accessing-fields:1}

\begin{lstlisting}
#eval intRing.addGrp.op 3 4     -- 7  (addition, via the nested CommGroup)
#eval intRing.mul 3 4            -- 12 (multiplication)
#eval intRing.one                 -- 1
#eval intRing.addGrp.toGroup.inv 5   -- -5  (additive inverse, via Group inside CommGroup)
\end{lstlisting}

\begin{mathreading}

The chained projections walk down the tower of \hyperref[sec:01-basics:04-terminology:1]{forgetful functors} \(\mathbf{Ring} \to \mathbf{Ab} \to \mathbf{Grp}\): \texttt{intRing.addGrp} applies \(\mathbf{Ring}\to\mathbf{Ab}\) (recovering \((\mathbb{Z},
+)\)), and \texttt{.toGroup} applies \(\mathbf{Ab}\to\mathbf{Grp}\). Thus, \texttt{intRing.addGrp.toGroup.inv 5} computes the additive inverse \(-5\) in the underlying group, while \texttt{intRing.mul}/\texttt{intRing.one} read off the multiplicative data \(\times\) and \(1\). Nested dot-access is just evaluating these structure-forgetting maps.

\end{mathreading}

\textbf{Mathlib equivalent.} There is no chain of \texttt{.addGrp.toGroup.inv}-style projections to write --- \texttt{+}/\texttt{*}/\texttt{0}/\texttt{1}/\texttt{-} already resolve to \texttt{Int}'s \texttt{CommRing} instance directly:

\begin{lstlisting}
#eval (3 : Int) + 4     -- 7
#eval (3 : Int) * 4      -- 12
#eval (1 : Int)          -- 1
#eval -(5 : Int)          -- -5
\end{lstlisting}

The book's nested projections walk down a tower of structures built by hand (\texttt{Ring → CommGroup → Group}). Mathlib's typeclass resolution walks down a similar tower of instances (\texttt{CommRing → Ring → ... → AddCommGroup}) automatically, so the notation never needs to name which layer it is coming from.
